% !TEX root = ../main.tex
\documentclass[../main.tex]{subfiles}
\begin{document}

\chapter{Introducción}
\label{chap:introduccion}

\section{Motivación}

La visión artificial embebida constituye uno de los campos de mayor crecimiento dentro de la ingeniería electrónica actual. 
La capacidad de adquirir, procesar y transmitir imágenes en tiempo real desde plataformas de bajo coste y reducidas dimensiones abre la puerta a multitud de aplicaciones industriales, médicas y de seguridad. Sin embargo, el diseño de estos sistemas presenta retos técnicos significativos: el ancho de banda necesario para transportar flujos de vídeo sin comprimir, la latencia introducida por las interfaces de comunicación y la correcta sincronización entre dominios de reloj heterogéneos son problemas que requieren soluciones cuidadosamente diseñadas a nivel de hardware.

En este contexto, las \ac{FPGA}s ofrecen ventajas únicas frente a los microcontroladores o los procesadores de propósito general: paralelismo inherente, latencia determinista y la posibilidad de implementar lógica de captura y transmisión de imagen sin necesidad de sistemas operativos ni \textit{drivers} complejos.

La Basys~3, basada en la familia Artix-7 de Xilinx, es una plataforma accesible y ampliamente utilizada en entornos académicos que, combinada con los periféricos adecuados, permite demostrar un sistema de visión embebido completo.

\section{Objetivos}

El presente \ac{TFM} tiene como objetivo principal el diseño, implementación y verificación de \ac{VICON}, un sistema de adquisición y transmisión de imagen configurable basado en \ac{FPGA}. Los objetivos específicos son:

\begin{itemize}
    \item Diseñar en \ac{VHDL} el hardware necesario para capturar imágenes del sensor CMOS MT9V111, incluyendo la interfaz de configuración I2C y el módulo de captura de trama.
    \item Implementar un controlador de alta velocidad para el chip FT232H de FTDI en modo \textit{Synchronous \ac{FIFO}}, capaz de transmitir el flujo de imagen hacia un PC y recibir comandos en sentido contrario.
    \item Gestionar correctamente los cruces de dominio de reloj (\ac{CDC}) entre los distintos dominios presentes en el diseño (reloj del sistema, reloj de píxel y reloj del FT232H).
    \item Verificar el diseño mediante simulación con metodología \ac{UVVM}, incluyendo agentes de simulación para el sensor I2C y el FT232H.
    \item Desarrollar una aplicación de visualización en PC con Qt y OpenCV que reciba el flujo de imagen y permita enviar comandos de control a la FPGA.
    \item Demostrar el funcionamiento del sistema completo en hardware real sobre la plataforma Basys~3.
\end{itemize}

\section{Estructura del documento}

El presente documento se organiza de la siguiente forma: 



El capítulo~\ref{chap:estado_del_arte} revisa el estado del arte en sistemas de visión embebida sobre FPGA, las interfaces de alta velocidad hacia PC y las metodologías de verificación de hardware.
El capítulo~\ref{chap:descripcion_sistema} describe la arquitectura global del sistema \ac{VICON} y el hardware utilizado.
El capítulo~\ref{chap:desarrollo} detalla el diseño a nivel \ac{RTL} de cada módulo, la estrategia de \ac{CDC} y la verificación mediante \ac{UVVM}.
El capítulo~\ref{chap:implementacion} presenta los resultados de síntesis e implementación en Vivado, el análisis de timing y las pruebas sobre hardware
real. 
El capítulo~\ref{chap:software} describe el software de visualización desarrollado en PC. 
Finalmente, el capítulo~\ref{chap:conclusiones} recoge las conclusiones del trabajo y las líneas de trabajo futuro.

\end{document}
